\documentclass[11pt,a4paper]{article}
\usepackage[margin=2.25cm,headheight=22pt]{geometry}
\usepackage[T1]{fontenc}
\usepackage[utf8]{inputenc}
\usepackage[spanish,es-nodecimaldot]{babel}
\usepackage{lmodern}
\usepackage{microtype}
\usepackage{xcolor}
\usepackage{graphicx}
\usepackage{tabularx}
\usepackage{array}
\usepackage{enumitem}
\usepackage{fancyhdr}
\usepackage{hyperref}
\usepackage{titlesec}

\definecolor{azul}{HTML}{18324B}
\definecolor{turquesa}{HTML}{087E83}
\definecolor{gris}{HTML}{52616B}
\definecolor{fondo}{HTML}{EEF4F6}
\definecolor{linea}{HTML}{CFDCE1}
\definecolor{pendiente}{HTML}{A25B00}

\hypersetup{
  colorlinks=true,
  linkcolor=turquesa,
  urlcolor=turquesa,
  pdftitle={Reporte de Hardening en Windows},
  pdfauthor={Juan Skobalski}
}

\pagestyle{fancy}
\fancyhf{}
\fancyhead[L]{\small\color{gris} Ciberseguridad | Unidad 2 | Ejercicio 2}
\fancyhead[R]{\small\color{gris} Juan Skobalski}
\fancyfoot[C]{\small\color{gris}\thepage}
\renewcommand{\headrulewidth}{0.4pt}
\renewcommand{\headrule}{\hbox to\headwidth{\color{linea}\leaders\hrule height \headrulewidth\hfill}}
\titleformat{\section}[block]{\Large\bfseries\color{azul}}{}{0pt}{}
\titleformat{\subsection}[block]{\large\bfseries\color{turquesa}}{}{0pt}{}
\titlespacing*{\section}{0pt}{0pt}{8pt}
\titlespacing*{\subsection}{0pt}{14pt}{6pt}
\setlist[itemize]{leftmargin=1.35em,itemsep=3pt,topsep=4pt}
\setlength{\parskip}{6pt}
\setlength{\parindent}{0pt}
\renewcommand{\arraystretch}{1.25}

% El documento solo se identifica como final cuando existen las tres capturas.
\newif\ifreportcomplete
\IfFileExists{capturas/01-usuario-estandar.png}{%
  \IfFileExists{capturas/02-windows-update.png}{%
    \IfFileExists{capturas/03-smb1-deshabilitado.png}{\reportcompletetrue}{\reportcompletefalse}%
  }{\reportcompletefalse}%
}{\reportcompletefalse}

\newcommand{\evidenceblock}[2]{%
  \begin{center}
  \IfFileExists{#1}{%
    \includegraphics[width=0.94\textwidth,height=8.0cm,keepaspectratio]{#1}\\[3pt]
    {\small\color{gris}#2}
  }{%
    \fcolorbox{linea}{fondo}{%
      \parbox[c][6.5cm][c]{0.88\textwidth}{%
        \centering
        {\Large\color{pendiente}\textbf{Evidencia pendiente}}\\[8pt]
        {\small Guardar la captura como \path{#1}}}}
  }
  \end{center}
}

\newcommand{\controlsummary}[3]{%
  \begin{tabularx}{\textwidth}{>{\raggedright\arraybackslash\bfseries\color{azul}}p{3.0cm} >{\raggedright\arraybackslash}X}
    Configuración & #1\\
    Objetivo & #2\\
    Verificación & #3\\
  \end{tabularx}
}

\begin{document}

% Carátula
\thispagestyle{empty}
\vspace*{1.0cm}
{\color{turquesa}\large C O D E R H O U S E}\par
\vspace{0.4cm}
{\color{linea}\rule{\textwidth}{1.2pt}}\par
\vspace{2.3cm}
{\color{gris}\Large CIBERSEGURIDAD}\par
\vspace{0.45cm}
{\color{azul}\Huge\bfseries Reporte de Hardening\\[5pt] en Windows}\par
\vspace{0.45cm}
{\color{turquesa}\LARGE Menor privilegio, actualizaciones y SMB 1.0}\par
\vspace{1.0cm}
{\color{gris}\large Unidad 2 | Ejercicio 2}\par

\vspace{0.8cm}
\ifreportcomplete
  \colorbox{fondo}{\parbox{\dimexpr\textwidth-2\fboxsep}{\centering\color{turquesa}\textbf{REPORTE CON EVIDENCIAS COMPLETAS}}}
\else
  \colorbox{fondo}{\parbox{\dimexpr\textwidth-2\fboxsep}{\centering\color{pendiente}\textbf{BORRADOR - EVIDENCIAS PENDIENTES}}}
\fi

\vfill
\noindent\colorbox{fondo}{%
  \begin{minipage}{\dimexpr\textwidth-2\fboxsep}
    \vspace{0.28cm}
    \noindent\makebox[3.1cm][l]{\textcolor{azul}{\textbf{Estudiante}}}Juan Skobalski\par
    \vspace{0.18cm}\color{linea}\hrule\color{black}\vspace{0.18cm}
    \noindent\makebox[3.1cm][l]{\textcolor{azul}{\textbf{Curso}}}Ciberseguridad\par
    \vspace{0.18cm}\color{linea}\hrule\color{black}\vspace{0.18cm}
    \noindent\makebox[3.1cm][l]{\textcolor{azul}{\textbf{Fecha}}}15 de octubre de 2026\par
    \vspace{0.28cm}
  \end{minipage}}
\vspace{0.75cm}
{\color{linea}\rule{\textwidth}{1.2pt}}\par
\vspace{0.25cm}
{\small\color{gris}Aplicación y documentación de controles básicos de endurecimiento en una máquina virtual de Windows.}

\newpage
\section*{Introducción}
\leavevmode\par

El \textbf{hardening} es el proceso de reducir la superficie de ataque de un sistema mediante configuraciones seguras, eliminación de funciones innecesarias, actualización del software y aplicación del principio de menor privilegio. Esta práctica se realiza en una máquina virtual de Windows para evitar cambios directos sobre el equipo anfitrión y facilitar la recuperación mediante instantáneas.

\subsection*{Objetivo y alcance}

El ejercicio documenta tres controles básicos:

\begin{enumerate}
  \item creación de una cuenta estándar separada de la cuenta administradora;
  \item ejecución de Windows Update para aplicar correcciones de seguridad;
  \item desactivación de SMB 1.0/CIFS cuando no existe una necesidad de compatibilidad.
\end{enumerate}

El reporte se limita a estos controles y no constituye una auditoría integral de Windows. No evalúa el firewall, el antivirus, el cifrado de disco ni las políticas corporativas del sistema.

\subsection*{Criterio de evaluación}

\begin{tabularx}{\textwidth}{>{\raggedright\arraybackslash\bfseries\color{azul}}p{3.2cm} >{\raggedright\arraybackslash}X}
  Menor privilegio & Los usuarios deben disponer únicamente de los permisos necesarios para su actividad cotidiana.\\
  Gestión de parches & El sistema debe recibir actualizaciones que corrijan vulnerabilidades conocidas.\\
  Superficie de ataque & Los protocolos y componentes heredados que no sean necesarios deben permanecer deshabilitados.\\
  Evidencia & Cada control debe estar respaldado por una captura legible de la VM de Windows.\\
\end{tabularx}

\subsection*{Resumen de controles}

\begin{tabularx}{\textwidth}{>{\raggedright\arraybackslash\bfseries\color{azul}}p{3.7cm} >{\raggedright\arraybackslash}X >{\raggedright\arraybackslash}X}
  Control & Riesgo reducido & Evidencia esperada\\
  Cuenta estándar & Abuso de privilegios administrativos. & Tipo de cuenta visible.\\
  Windows Update & Explotación de fallas ya corregidas. & Estado de actualización.\\
  SMB 1.0 deshabilitado & Exposición de un protocolo heredado. & Casilla SMB 1.0/CIFS vacía.\\
\end{tabularx}

\vfill
\noindent\colorbox{fondo}{\parbox{\dimexpr\textwidth-2\fboxsep}{\textcolor{azul}{\textbf{Precaución:}} antes de modificar cuentas o características de Windows, se debe crear una instantánea y comprobar que la contraseña de la cuenta administradora funciona.}}

\newpage
\section*{Aplicación de los controles}

\subsection*{1. Cuenta de uso diario sin privilegios administrativos}

\controlsummary
{Cuenta local \texttt{standard} configurada como \textbf{Usuario estándar}.}
{Aplicar menor privilegio para que las tareas cotidianas no se ejecuten con permisos administrativos.}
{En \textit{Configuración > Cuentas > Otros usuarios}, comprobar el nombre de la cuenta y su tipo. La captura no debe mostrar contraseñas ni datos de recuperación.}

\evidenceblock{capturas/01-usuario-estandar.png}{Figura 1. Cuenta de uso diario configurada como usuario estándar.}

\newpage
\section*{Actualización y gestión de parches}

\subsection*{2. Actualización del sistema operativo}

\controlsummary
{Búsqueda e instalación de actualizaciones desde Windows Update.}
{Corregir vulnerabilidades conocidas y mantener vigentes los componentes de seguridad del sistema.}
{Comprobar que Windows Update esté buscando actualizaciones o informe que el sistema se encuentra actualizado. Reiniciar si el sistema lo solicita.}

\evidenceblock{capturas/02-windows-update.png}{Figura 2. Estado de Windows Update dentro de la máquina virtual.}

\newpage
\section*{Reducción de la superficie de ataque}

\subsection*{3. Desactivación de SMB 1.0/CIFS}

\controlsummary
{Compatibilidad con el protocolo SMB 1.0/CIFS deshabilitada.}
{Evitar la exposición innecesaria de un protocolo heredado para compartir archivos.}
{Abrir \textit{Activar o desactivar las características de Windows} y confirmar que la casilla de SMB 1.0/CIFS esté vacía.}

\evidenceblock{capturas/03-smb1-deshabilitado.png}{Figura 3. Característica SMB 1.0/CIFS deshabilitada.}

\subsection*{Resultado y conclusión}

\ifreportcomplete
Las evidencias incorporadas documentan la separación entre el uso cotidiano y la administración, el funcionamiento del mecanismo de actualización y la desactivación de un protocolo heredado. En conjunto, estos controles reducen privilegios, corrigen vulnerabilidades conocidas y eliminan una función innecesaria de la superficie de ataque.
\else
La estructura del reporte y el análisis técnico están preparados, pero todavía faltan una o más capturas de la máquina virtual. El documento no debe presentarse como evidencia final hasta reemplazar todos los recuadros pendientes y verificar que las imágenes permitan confirmar cada control.
\fi

El hardening debe tratarse como un proceso continuo. Después de instalar aplicaciones, cambiar servicios o aplicar nuevas actualizaciones, corresponde revisar nuevamente la configuración y conservar un punto de restauración funcional.

\subsection*{Checklist de entrega}

\begin{itemize}
  \item La cuenta de uso diario figura como usuario estándar.
  \item Windows Update aparece funcionando o con el sistema actualizado.
  \item SMB 1.0/CIFS aparece deshabilitado.
  \item Las capturas son legibles y no muestran información sensible.
  \item El PDF se entrega como \texttt{Reporte\_Hardening\_Skobalski\_Juan.pdf}.
\end{itemize}

\end{document}
